\section{Exact finite reporters}
\label{sec:reporter}

We first state the cancellation mechanism from
\cite[Sections 1--2]{OpenAI2026}, and then derive a closed finite-law
formula that will be used for the explicit witness.

\subsection{The reporting rule}

Let \(U,Z_1,\ldots,Z_m\) be independent copies of a centered random
variable with finite support \(\mathcal S\) and variance \(\sigma^2\).
Choose leaf tests \(I_i\subseteq\mathcal S\) and a deterministic selector
\(\iota:\mathcal S\to[m]\). Define
\[
 B=\{Z_j\in I_j\ \hbox{for all }j\},\qquad
 P=\{Z_j\in I_j\ \hbox{for all }j\ne\iota(U)\},\qquad
 A=P\setminus B.
\]
The report \(\mathcal W\) is the following tagged observation.
\begin{center}
\begin{tabular}{@{}ll@{}}
\toprule
Event & Reported information\\
\midrule
\(A\) & The single tag \(A\)\\
\(B\) & The tag \(B\) and all \(m\) leaf values\\
\(P^c\) & The ordinary tag, \(U\), \(\iota(U)\), and all unselected leaves\\
\bottomrule
\end{tabular}
\end{center}
Put \(L=U+\sum_j Z_j\). The exceptional indicator satisfies the
pointwise identity
\begin{equation}\label{eq:reporter-cancellation}
 \1_A=
 \sum_{i=1}^m\1_{\{\iota(U)=i\}}
       \prod_{j\ne i}\1_{\{Z_j\in I_j\}}
 -\prod_{j=1}^m\1_{\{Z_j\in I_j\}}.
\end{equation}
Each summand involves at most \(m\) of the \(m+1\) inputs.

\begin{proposition}[Degree cost and exact exceptional gain]
\label{prop:reporter}
Suppose each scalar input is a function of an independent cube
observation of cell degree at most \(d\). Then every cell of
\(\mathcal W\) has degree at most \(md\). If \(\Pbb(A)>0\), then
\begin{equation}\label{eq:optimal-exceptional}
 \Var(\E[L\mid\mathcal W])
 =m\sigma^2+\Pbb(A)\bigl(\sigma^2-\Var(L\mid A)\bigr).
\end{equation}
For any prescribed exceptional predictor \(b\),
\begin{equation}\label{eq:reporter-predictor}
 \Var(\E[L\mid\mathcal W])
 \ge m\sigma^2+
 \E\bigl[\1_A\{\sigma^2-(L-b)^2\}\bigr].
\end{equation}
\end{proposition}

\begin{proof}
Equation~\eqref{eq:reporter-cancellation} and
Lemma~\ref{lem:observation-calculus} handle the exceptional cell.
A \(B\)-cell fixes all leaves and ignores \(U\). An ordinary cell fixes
\(U\) and the unselected leaves, at least one of which fails its test.
Those fixed values already force \(P^c\), so the selected leaf remains
unrestricted. These indicators involve at most \(m\) inputs.

On each nonexceptional cell, precisely one independent input with its
original law is unreported. Thus the conditional variance of \(L\) is
\(\sigma^2\) there. On \(A\) it is \(\Var(L\mid A)\).
Subtracting the averaged residual variance from
\(\Var L=(m+1)\sigma^2\) gives
\eqref{eq:optimal-exceptional}. Replacing the optimal mean on \(A\)
by \(b\) only increases the mean squared error and proves the last
inequality.
\end{proof}

The same predictor calculation holds for real input laws, with finite
interval tests and a measurable selector. Only the assertion about
finite cell indicators needs the finite-alphabet restriction. This
distinction will be used for the Gaussian model.

\subsection{An exact formula for singleton tests}

Write \(p(t)=\Pbb(U=t)\). Suppose \(I_i=\{a_i\}\), with
\(p(a_i)>0\), and set
\begin{align}
 Q&=\prod_{i=1}^m p(a_i),&
 a_\Sigma&=\sum_{i=1}^m a_i,&
 w_t&=\frac{p(t)}{p(a_{\iota(t)})},\notag\\
 W&=\sum_{t\in\mathcal S}w_t,&
 K_0&=W-1,&
 R_1&=\sum_{t\in\mathcal S}w_t(t-a_{\iota(t)}),\notag\\
 H_2&=\sum_{t\in\mathcal S}w_t(t-a_{\iota(t)})^2.&&&
 \label{eq:singleton-parameters}
\end{align}
The notation \(K_0\) is local to this finite formula; it is not the
fourth-moment bound used later.

\begin{theorem}[Finite selector formula and score law]
\label{thm:finite-selector}
Assume \(K_0>0\). The exceptional report has probability \(QK_0\) and
conditional mean
\[
 b_*=\E[L\mid A]=a_\Sigma+\frac{R_1}{K_0}.
\]
For \(T=\E[L\mid\mathcal W]\),
\begin{equation}\label{eq:finite-retention}
 \E T^2=m\sigma^2+
 Q\left(\frac{R_1^2}{K_0}-H_2\right).
\end{equation}
In particular, the report has positive variance advantage if and only
if \(R_1^2>K_0H_2\). More strongly,
\begin{equation}\label{eq:finite-score-law}
 \Law(T)=\Law(U_1+\cdots+U_m)
 -Q\sum_{t\in\mathcal S}w_t\delta_{a_\Sigma+t-a_{\iota(t)}}
 +Q\delta_{a_\Sigma}+QK_0\delta_{b_*}.
\end{equation}
The right side is an identity of signed measures whose total is the
probability law on the left.
\end{theorem}

\begin{proof}
Given \(U=t\), event \(P\) has probability
\(Q/p(a_{\iota(t)})\). Thus
\[
 \Pbb(P)=QW,\quad \Pbb(B)=Q,\quad \Pbb(A)=QK_0.
\]
On \(P\), the unselected leaf sum equals
\(a_\Sigma-a_{\iota(t)}\), and the selected input is unrestricted.
Therefore
\[
 \E[L\1_P]=Q\sum_t w_t(a_\Sigma+t-a_{\iota(t)}).
\]
On \(B\), all leaves are fixed while \(U\) is centered, so
\(\E[L\1_B]=Qa_\Sigma\). Subtraction and division by \(QK_0\)
give \(b_*\).

Define the baseline random variable
\[
 T_0=U+\sum_{j\ne\iota(U)} Z_j.
\]
Conditionally on \(U=t\), the remaining sum always has the law of
\(m-1\) independent inputs. It follows that
\(\Law(T_0)=\Law(U_1+\cdots+U_m)\), despite the
dependence of the omitted index on \(U\).
The true score agrees with \(T_0\) on \(P^c\).
The restriction of the law of \(T_0\) to \(P\) is
\[
 Q\sum_t w_t\delta_{a_\Sigma+t-a_{\iota(t)}}.
\]
On \(B\) the true score is \(a_\Sigma\), and on \(A\) it is \(b_*\).
Replacing these restrictions gives \eqref{eq:finite-score-law}.

The first moment of the correction is zero. For the second moment,
the terms involving \(a_\Sigma^2\) cancel because
\(-W+1+K_0=0\); the linear terms cancel because
\(K_0(b_*-a_\Sigma)=R_1\). The remaining correction is precisely
\(Q(R_1^2/K_0-H_2)\), proving \eqref{eq:finite-retention}.
\end{proof}

This formula turns a probabilistic design problem into a finite
rational inequality. It also provides the entire score distribution,
so higher moments and readout behavior can be checked without
enumerating every bit string.
